\documentclass[11pt,a4paper]{article}

\usepackage[utf8]{inputenc}
\usepackage[T1]{fontenc}
\usepackage[margin=2.2cm]{geometry}
\usepackage{parskip}
\usepackage{titlesec}
\usepackage{enumitem}
\usepackage{booktabs}
\usepackage{xcolor}
\usepackage{tcolorbox}
\usepackage{hyperref}
\usepackage{fancyhdr}
\usepackage{listings}
\usepackage{longtable}
\usepackage{array}

% Long unbreakable identifiers are common here; let TeX stretch spaces instead of
% running into the margin.
\setlength{\emergencystretch}{3em}

\tcbuselibrary{skins,breakable}

\definecolor{accent}{RGB}{30,80,150}
\definecolor{softgray}{RGB}{245,246,248}
\definecolor{codebg}{RGB}{248,248,250}
\definecolor{warn}{RGB}{160,60,40}
\definecolor{good}{RGB}{30,120,70}

\titleformat{\section}{\Large\bfseries\color{accent}}{\thesection}{0.8em}{}
\titleformat{\subsection}{\large\bfseries}{\thesubsection}{0.6em}{}
\titleformat{\subsubsection}{\normalsize\bfseries\itshape}{\thesubsubsection}{0.5em}{}

\hypersetup{
  colorlinks=true, linkcolor=accent, urlcolor=accent,
  pdftitle={AgentIA --- Plan de Pruebas y Resultados de Aseguramiento de Calidad},
  pdfauthor={Equipo de Ingenieria AgentIA}
}

\pagestyle{fancy}
\fancyhf{}
\fancyhead[L]{\small\color{accent}AgentIA --- Plan de Pruebas y Resultados QE}
\fancyhead[R]{\small\color{accent}\thepage}
\renewcommand{\headrulewidth}{0.4pt}

\lstset{
  basicstyle=\ttfamily\small,
  backgroundcolor=\color{codebg},
  frame=single,
  rulecolor=\color{softgray},
  breaklines=true,
  columns=fullflexible,
  keepspaces=true,
  showstringspaces=false,
}

\newtcolorbox{finding}[1][]{
  colback=softgray, colframe=accent, boxrule=0.6pt, arc=2pt,
  left=6pt, right=6pt, top=5pt, bottom=5pt, breakable, #1
}

% Underscores in identifiers carry no break opportunity in \ttfamily, so long test names
% run into the margin. Redefining \_ locally restores one.
\newcommand{\code}[1]{{\ttfamily\small\def\_{\textunderscore\allowbreak}#1}}
\newcommand{\codep}[1]{\code{#1}}

\title{\vspace{-1.5cm}\textbf{Plan de Pruebas y Resultados}\\[2pt]
\large Campaña de Aseguramiento de Calidad (QE)\\[6pt]
\normalsize Microservice Code Studio --- AgentIA}
\author{Equipo de Ingeniería AgentIA}
\date{29 de septiembre de 2026}

\begin{document}
\maketitle
\vspace{-0.6cm}

\begin{finding}
\textbf{Para revisión del Project Manager.} Este documento describe qué se probó, cómo,
qué se encontró y qué falta. Las cifras de la Sección~\ref{sec:resultados} se obtuvieron
ejecutando la suite completa dos veces sobre el mismo entorno: una sobre el código anterior
(punto de partida) y otra sobre el código actual. El comando exacto de reproducción está en
la Sección~\ref{sec:reproduccion}. Todos los registros crudos quedan versionados en
\code{reports/qe/}.
\end{finding}

\tableofcontents
\newpage

% ===========================================================================
\section{Resumen ejecutivo}
\label{sec:resumen}

La campaña tuvo dos fases. La primera cubrió los tres módulos con mayor riesgo de producto
según el inventario de la campaña anterior (\code{docs/qe\_test\_matrix.md}); la segunda cerró
el motor del Auto-Pilot y añadió pruebas de flujo del frontend en la frontera HTTP. El
criterio no fue ``el número más bajo'', sino \emph{qué decisión toma cada módulo y cuánto
duele si la toma mal}.

\begin{center}
\begin{tabular}{lrr}
\toprule
\textbf{Módulo} & \textbf{Antes} & \textbf{Después} \\
\midrule
\code{services/docker\_service.py}      & 52{,}7\,\% & \textbf{100{,}0\,\%} \\
\code{api/routes\_session.py}           & 71{,}3\,\% & \textbf{100{,}0\,\%} \\
\code{services/lifecycle\_service.py}   & 76{,}4\,\% & \textbf{100{,}0\,\%} \\
\code{services/pipeline\_runner.py}     & 82{,}4\,\% & \textbf{100{,}0\,\%} \\
\midrule
Sentencias sin cubrir                   & 994        & \textbf{727} \\
Pruebas de backend                      & 589 correctas & \textbf{825 correctas} \\
Pruebas de frontend                     & 64 (10 archivos) & \textbf{73 (11 archivos)} \\
\bottomrule
\end{tabular}
\end{center}

\noindent
La campaña añadió \textbf{186 pruebas automatizadas} (177 de backend, 9 de frontend) y
encontró \textbf{cinco defectos reales}. Ninguno era visible leyendo el código: todos
aparecieron porque una prueba falló, se colgó, o ejecutó por primera vez una rama que nunca
se había ejecutado.

\begin{itemize}[leftmargin=1.4em]
\item \textbf{D-03 (crítico).} El espejo de telemetría de costos, documentado como
      ``best-effort, fuera de toda ruta crítica'', \emph{bloqueaba} la ruta crítica. Un
      destino inalcanzable detenía cada llamada al modelo durante minutos, y detenía la
      suite completa de forma indefinida. Con el defecto presente, la cobertura del backend
      no podía medirse en absoluto.
\item \textbf{D-05 (alto).} El respaldo de síntesis de esquema SQL --- añadido precisamente
      para no entregar un esquema equivocado --- llamaba a un método que no existe, dejaba un
      \code{schema.sql} \emph{vacío} y tumbaba la ejecución completa.
\item \textbf{D-02 (alto).} Las fases 5 y 6 mostraban el estado \code{BLOCKED} con
      \code{blockingReason: null}: el operador veía que el pipeline se había detenido, pero
      no por qué. Afecta directamente al Principio V de la constitución.
\item \textbf{D-04 (alto).} La decisión de modo y la carga del conjunto de instrucciones
      ocurrían \emph{fuera} del manejador de errores del pipeline, de modo que un fallo ahí
      escapaba del hilo demonio y dejaba la sesión en \code{RUNNING} para siempre.
\item \textbf{D-01 (medio).} Dos funciones del ciclo de vida estaban definidas dos veces en
      el mismo archivo. La primera definición era código muerto e imposible de cubrir, y su
      contrato \emph{difería} del de la segunda.
\end{itemize}

\noindent
Los cinco están corregidos y fijados con pruebas de regresión. La Sección~\ref{sec:defectos}
documenta causa raíz y evidencia de cada uno.

% ===========================================================================
\section{Objetivo y alcance}

\subsection{Objetivo}

Incorporar pruebas automatizadas ejecutables y deterministas sobre las superficies del
producto con mayor cobertura pendiente, y entregar evidencia verificable de su resultado
para la revisión del Project Manager.

\subsection{Alcance acordado}

La campaña cubre \textbf{tres módulos de backend}, elegidos por riesgo:

\begin{enumerate}[leftmargin=1.6em]
\item \code{services/docker\_service.py} --- la ruta de despliegue local que la interfaz de
      demostración opera directamente (40{,}4\,\% de cobertura según la campaña anterior).
\item \code{api/routes\_session.py} --- la superficie principal del producto: creación,
      listado, lectura, cancelación de sesiones, el flujo SSE del monitor y los caminos de
      desenlace de la generación (69{,}9\,\%).
\item \code{services/lifecycle\_service.py} --- la máquina de estados de fases: qué se puede
      entrar, qué porcentaje está completo y qué se recomienda a continuación (76{,}4\,\%).
\end{enumerate}

\subsection{Fuera de alcance}

\begin{itemize}[leftmargin=1.4em]
\item Pruebas de aceptación independientes contra un microservicio generado y desplegado
      (la limitación de la Sección~\ref{sec:limitaciones}); requieren Docker y una
      compilación real.
\item Pruebas de navegador de extremo a extremo (Playwright). El frontend \emph{sí} recibió
      pruebas de flujo, pero en jsdom y contra un backend simulado, no en un navegador real
      contra servidores vivos.
\item \code{ssl\_compat.py} y otras rutas de valor bajo.
\end{itemize}

% ===========================================================================
\section{Estrategia de pruebas}

\subsection{Enfoque basado en riesgo}

Cada módulo se abordó preguntando primero \emph{qué afirmación hace} y \emph{qué pasa si la
hace mal}, y sólo después qué líneas faltaban. Las afirmaciones que se fijaron son:

\begin{itemize}[leftmargin=1.4em]
\item \code{docker\_service}: \code{HEALTHY} es una \textbf{afirmación de verificación}. Sólo
      es honesta si un \code{/actuator/health} real respondió \code{UP}. Que
      \code{docker compose up} salga con código 0 significa ``los contenedores arrancaron'',
      nunca ``el servicio funciona''.
\item \code{routes\_session}: los caminos de desenlace deben decir la verdad. Una sesión que
      llegó a \code{VERIFIED} sobre una verificación \emph{sustituida} debe declararlo
      (\code{verificationFallbackUsed}, feature 012); una sesión bloqueada debe registrarse
      igual que una completada (feature 015); un fallo dentro del grafo debe terminar la
      sesión, no dejarla en \code{RUNNING} para siempre.
\item \code{lifecycle\_service}: los prerrequisitos de fase son una regla de negocio. El
      límite de 3 auto-reparaciones (Principio V) debe cerrar la puerta, no sólo pintarse.
\end{itemize}

\subsection{Sustitución de fronteras}

Ninguna prueba nueva inicia un contenedor, abre un socket ni compila Java. Las fronteras
externas se sustituyen siempre en el mismo punto:

\begin{center}
\begin{tabular}{lll}
\toprule
\textbf{Frontera} & \textbf{Sustituto} & \textbf{Motivo} \\
\midrule
\code{subprocess.Popen/run} & proceso falso con salida guionada & sin Docker \\
\code{requests.get}        & respuesta falsa o excepción     & sin red \\
\code{threading.Thread}    & hilo en línea                   & determinismo \\
\code{generation\_graph}   & grafo guionado                  & sin LLM \\
\code{audit\_workspace}    & auditoría falsa                 & sin SAST \\
\bottomrule
\end{tabular}
\end{center}

\noindent
El hilo del \code{docker compose} merece mención aparte. \code{deploy\_local} devuelve el
control en cuanto lanza el hilo trabajador; una prueba que esperase a que terminara sería
lenta en una máquina descargada y fallida en una cargada. Sustituyendo \code{Thread} por una
implementación que ejecuta el objetivo en línea, la prueba observa un despliegue
\emph{terminado} sin dormir ni sondear.

\subsection{Hermeticidad}

La suite quedó aislada de la red. Las tres tiendas de estado a nivel de módulo
(\code{\_active\_deployments}, \code{SESSION\_EVENT\_HISTORY}, \code{SESSION\_EVENT\_SUBSCRIBERS})
se aíslan por prueba, porque son globales al proceso por diseño y sin aislamiento una prueba
devuelve la sesión de otra. El destino de telemetría se redirige a un almacén local temporal
(Sección~\ref{sec:d03}).

% ===========================================================================
\section{Inventario de pruebas automatizadas}
\label{sec:inventario}

\begin{center}
\footnotesize
\begin{longtable}{>{\raggedright\arraybackslash}p{4.6cm}r >{\raggedright\arraybackslash}p{9.0cm}}
\toprule
\textbf{Archivo} & \textbf{Pruebas} & \textbf{Superficie que cubre} \\
\midrule
\endfirsthead
\toprule
\textbf{Archivo} & \textbf{Pruebas} & \textbf{Superficie que cubre} \\
\midrule
\endhead
\code{test\_qe\_deploy\_docker.py} & 33 &
Detección del demonio Docker; \code{deploy\_local} (éxito, sin demonio, \code{compose} con
error, smoke test fallido, excepción del trabajador, \code{rebuild}); \code{run\_smoke\_test}
(UP, estado degradado, cuerpo no JSON, reintentos, éxito tardío);
\code{get\_deployment\_status}; \code{stream\_logs}; \code{stop\_deployment}. \\
\addlinespace
\code{test\_qe\_session\_surface.py} & 50 &
Lectura y persistencia de \code{verificationFallbackUsed}; difusión de eventos SSE y
suscriptores caídos; \code{stream\_session\_events} (replay, evento en vivo, keep-alive,
desconexión); \code{list\_sessions} (promoción a COMPLETED, fallo tolerado);
\code{quick\_start\_session}; cancelación; y los cuatro desenlaces de
\code{execute\_generation\_pipeline}. \\
\addlinespace
\code{test\_qe\_lifecycle\_rules.py} & 46 &
Detección de las 7 fases; porcentaje de avance; prerrequisitos; límite de 3
auto-reparaciones; Quality Gate bloqueado y auditoría que falla; fases obsoletas
(marcado, acumulación, limpieza, JSON corrupto); degradación de valores de enumeración
desconocidos; estado del pipeline; \code{get\_project\_overview}. \\
\addlinespace
\code{test\_cost\_mlflow\_mirror.py} & +2 &
Cota de transporte del espejo de telemetría (sin reintentos, tiempo de espera acotado) y
respeto de una configuración explícita del operador. \\
\addlinespace
\code{test\_qe\_pipeline\_autopilot.py} & 46 &
Pausa cooperativa en cada frontera de paso; cancelar frente a pausar; los puntos de entrada
\code{pause}/\code{resume}/\code{cancel} y la guarda de concurrencia de \code{run\_pipeline};
\code{stream\_pipeline\_events} completo (conexión, progreso, heartbeats, cierre); el camino
de agotamiento; el respaldo de esquema y el de arquitectura; y la propagación de errores. \\
\addlinespace
\code{workflow\_journeys.test.tsx} (frontend) & 9 &
Flujos de usuario contra un backend simulado en la frontera HTTP con MSW: inicio rápido y
lectura de vuelta, ciclo de control del Auto-Pilot, resumen y ciclo de vida, artefactos,
credencial efímera, propagación de fallos, y el flujo en la interfaz incluido el informe
honesto del smoke test. \\
\bottomrule
\end{longtable}
\end{center}

\noindent
Total: \textbf{186 pruebas nuevas} (177 de backend, 9 de frontend). Las campañas anteriores
(\code{test\_qe\_api\_surface.py}, \code{test\_qe\_environment\_repair.py},
\code{test\_qe\_publish\_git.py}, \code{test\_qe\_schema\_consistency.py}) permanecen y siguen
pasando: 67 pruebas más.

% ===========================================================================
\section{Matriz de trazabilidad}
\label{sec:trazabilidad}

\begin{center}
\footnotesize
\begin{longtable}{>{\raggedright\arraybackslash}p{3.6cm} >{\raggedright\arraybackslash}p{11.8cm}}
\toprule
\textbf{Requisito / Principio} & \textbf{Afirmación verificada y prueba que la fija} \\
\midrule
\endfirsthead
\toprule
\textbf{Requisito / Principio} & \textbf{Afirmación verificada y prueba que la fija} \\
\midrule
\endhead
Principio V --- auto-reparación acotada a 3 intentos &
Con 3 intentos la fase 5 queda \code{BLOCKED} y \code{canAdvance} es falso; con 2 no. La
razón del bloqueo llega al cliente.
\newline\code{test\_a\_blocked\_repair\_loop\_blocks\_phase\_five...},
\code{test\_two\_failed\_repairs\_do\_not\_block\_phase\_five},
\code{test\_a\_repair\_beyond\_the\_limit\_is\_reported\_as\_blocked} \\
\addlinespace
Feature 012 (FR-005) --- la verificación sustituida debe declararse &
El indicador se lee del registro persistido (sobrevive a un reinicio), se escribe en el
evento terminal y es una medición, no una constante.
\newline\code{test\_reading\_a\_session\_surfaces\_the\_persisted\_fallback\_flag},
\code{test\_a\_completed\_run\_records\_verified\_and\_reports\_the\_substituted\_verification},
\code{test\_a\_completed\_run\_with\_a\_real\_verification\_reports\_no\_fallback} \\
\addlinespace
Feature 015 (FR-001) --- el diagnóstico se registra en todo camino terminal &
Tanto el camino \code{COMPLETED} como el \code{BLOCKED} invocan al escritor único.
\newline\code{test\_a\_completed\_run\_records\_verified\_and\_reports\_the\_substituted\_verification},
\code{test\_a\_blocked\_run\_is\_persisted\_and\_diagnosed} \\
\addlinespace
Honestidad de la verificación (feature 012) &
\code{compose} con código 0 y health degradado produce \code{FAILED}, nunca \code{HEALTHY}.
\newline\code{test\_a\_container\_that\_never\_becomes\_healthy\_is\_failed\_not\_healthy},
\code{test\_a\_missing\_health\_field\_is\_not\_treated\_as\_up} \\
\addlinespace
Principio VI --- cero secretos, cero red en pruebas &
La suite no realiza ninguna conexión saliente; el token de Git no sobrevive a la publicación.
\newline\code{no\_outbound\_telemetry} (conftest), \code{test\_qe\_publish\_git.py} \\
\addlinespace
Principio III --- sobre de error uniforme &
Los manejadores devuelven \code{status}/\code{message}/\code{details}.
\newline\code{test\_qe\_api\_surface.py} (campaña anterior) \\
\addlinespace
Honestidad de la verificación, de extremo a extremo &
Cuando el endpoint de salud falla, la interfaz muestra ``Smoke test NO ejecutado''; nunca un
éxito que nadie observó. Un fallo del backend se propaga en lugar de inventarse.
\newline\code{workflow\_journeys.test.tsx} \\
\bottomrule
\end{longtable}
\end{center}

% ===========================================================================
\section{Defectos encontrados y corregidos}
\label{sec:defectos}

Los tres defectos siguientes fueron encontrados \emph{por las pruebas}, no por lectura del
código. Se documentan con severidad, evidencia, causa raíz y prueba de regresión.

\subsection*{D-03 --- El espejo de telemetría ``best-effort'' bloqueaba la ruta crítica}
\label{sec:d03}

\begin{finding}
\textbf{Severidad: crítica.} Afecta a producción y al arnés de pruebas.
\end{finding}

\paragraph{Evidencia.} La suite completa no terminaba. Se detenía indefinidamente en
\code{test\_cost\_recording.py::}\allowbreak\code{test\_a\_call\_through\_the\_stage\_boundary\_is\_recorded\_end\_to\_end}.
La traza del proceso colgado, obtenida con \code{faulthandler}, muestra la pila exacta:

\begin{lstlisting}
urllib3/util/retry.py            _sleep_backoff
requests/adapters.py             send
mlflow/tracking/fluent.py        start_run
app/cost/mlflow_sink.py:49       _record
app/cost/recording.py:148        _record_call
app/orchestrator/stages/runner.py:994  _run_model_stage
backend/tests/test_cost_recording.py:116
\end{lstlisting}

\paragraph{Causa raíz.} \code{mlflow\_sink.py} se declara ``deliberadamente fuera de toda
ruta crítica'' y sostiene que un destino inalcanzable ``simplemente no hace nada''. El
transporte no se comportaba así. Los valores por defecto de MLflow son
\code{MLFLOW\_HTTP\_REQUEST\_TIMEOUT = 120} segundos y
\code{MLFLOW\_HTTP\_REQUEST\_MAX\_RETRIES = 7} con retroceso exponencial
(factor 2); la propia documentación de la librería calcula unos 4 minutos para esas 7
tentativas. El destino por defecto es \code{http://localhost:5000}, que en este entorno no
existe.

Y el llamador \emph{sí} está en la ruta crítica: \code{mirror\_call\_record} se invoca desde
\code{RecordingChatClient}\allowbreak\code{.invoke}, que envuelve \textbf{cada} llamada al
modelo. Por tanto,
una caída de telemetría detenía cada sesión de generación durante minutos. El módulo no era
ni ``fuera de ruta'' ni silencioso.

\paragraph{Corrección.} Se añadió \code{\_bound\_the\_transport()}, que fija las dos únicas
variables que MLflow consulta en cada petición
(\code{MLFLOW\_HTTP\_REQUEST\_TIMEOUT = 5},
\code{MLFLOW\_HTTP\_REQUEST\_MAX\_RETRIES = 0}) mediante \code{setdefault}. Se usa
\code{setdefault} y no asignación directa para que un despliegue con un servidor MLflow sano
pueda pedir reintentos largos sin que el espejo se los quite.

\paragraph{Corrección del arnés.} El destino se redirige a un almacén local temporal en
\code{conftest.py}. Se redirige en lugar de anularse porque
\code{test\_the\_local\_store\_is\_written\_even\_when\_the\_mirror\_is\_absent} afirma que el
espejo \emph{sí} fue consultado; anularlo habría vuelto vacía esa prueba. La dirección se
parchea sobre la configuración y no sobre el ayudante \code{\_tracking\_uri}, para que el
ayudante siga bajo prueba.

\paragraph{Resultado.} La prueba que colgaba indefinidamente ahora tarda 1{,}45 s. La suite
completa termina. Se añadieron dos pruebas de regresión que fijan la cota de transporte y el
respeto de la configuración explícita.

\subsection*{D-05 --- El respaldo del esquema llamaba a un método inexistente}

\begin{finding}
\textbf{Severidad: alta.} Rompía la ejecución completa y dejaba un artefacto vacío.
\end{finding}

\paragraph{Evidencia.} La primera prueba que ejecutó la rama de respaldo del esquema falló
con \code{AttributeError}. El archivo \code{schema.sql} existía, pero estaba \emph{vacío}.
La ejecución terminaba en \code{FAILED}.

\paragraph{Causa raíz.} \code{schema\_sql\_from\_draft} es una función de módulo, no un
método del objeto \code{model\_sql\_service}. El respaldo --- introducido por el commit
\code{37b5853} justamente para no entregar un esquema equivocado --- la invocaba como
\code{model\_sql\_service.schema\_sql\_from\_draft(draft)}, de modo que lanzaba
\code{AttributeError}. Y como \code{open(..., "w")} trunca el archivo \emph{antes} de
evaluar su argumento, el resultado era un \code{schema.sql} vacío además de la caída.

Nadie lo detectó porque la rama nunca se ejecutó: estaba en el conjunto de líneas sin cubrir.
La función sí estaba probada de forma aislada (en \code{test\_qe\_schema\_consistency.py}),
que es precisamente por lo que la prueba pasaba mientras el punto de llamada estaba roto.

\paragraph{Corrección.} Se importa la función correcta y el contenido se calcula antes de
abrir el archivo, de modo que un fallo ya no puede truncar el artefacto.

\paragraph{Pruebas de regresión.} La prueba fuerza el fallo de la síntesis y exige que la
ejecución \emph{complete}, que el esquema no esté vacío y que describa las entidades de
\emph{este} blueprint.

\subsection*{D-04 --- El fallo más ruidoso escapaba del manejador y no lo veía nadie}

\begin{finding}
\textbf{Severidad: alta.} Dejaba la sesión en \code{RUNNING} indefinidamente.
\end{finding}

\paragraph{Evidencia.} La prueba del modo MODEL con un conjunto de instrucciones ilegible
esperaba una sesión marcada como bloqueada con la causa registrada. En su lugar, la
excepción escapó de \code{\_execute\_pipeline\_steps}.

\paragraph{Causa raíz.} El \code{try} del pipeline empezaba \emph{después} de la selección de
modo y de la carga del conjunto de instrucciones. El código re-lanza deliberadamente en modo
MODEL para fallar con fuerza (invariante 11) --- y eso es correcto ---, pero al estar fuera
del manejador la excepción moría con el hilo demonio. La fila de sesión se quedaba en
\code{RUNNING} y el operador miraba un pipeline que ya se había detenido. Lo mismo aplicaba a
una clave de API ambigua, que \code{detect\_provider} rechaza a propósito.

\paragraph{Corrección.} El manejador ahora envuelve también la decisión de modo y la carga de
instrucciones. Fallar con fuerza sigue siendo el comportamiento; fallar \emph{en silencio
para el operador} no.

\paragraph{Pruebas de regresión.} Dos pruebas: el conjunto de instrucciones ilegible en modo
MODEL y la clave ambigua. Ambas exigen estado \code{FAILED} y la causa en
\code{error\_message}.

\subsection*{D-02 --- La razón del bloqueo se calculaba y se descartaba}

\begin{finding}
\textbf{Severidad: alta.} Afecta al Principio V y a la experiencia del operador.
\end{finding}

\paragraph{Evidencia.} La prueba del límite de 3 auto-reparaciones afirmaba que el texto de
la razón del bloqueo llegara al cliente. Falló: \code{state.phases[4].blocking\_reason} era
\code{None}. La fase se reportaba correctamente como \code{BLOCKED}, pero sin causa.

\paragraph{Causa raíz.} Ambas rutas de bloqueo asignaban una variable local
\code{blocked\_reason}, mientras que el objeto \code{PhaseState} se construye a partir de
\code{reason}. \code{blocked\_reason} se escribía tres veces (inicialización y las dos rutas
de bloqueo) y \textbf{nunca se leía}. El resultado en la API era
\code{blockingReason: null} para la fase 5 detenida por el límite constitucional y para la
fase 6 detenida por un Quality Gate bloqueado.

\paragraph{Corrección.} Ambas rutas asignan ahora \code{reason}, que es el campo que se
serializa. Se eliminó la variable muerta.

\paragraph{Pruebas de regresión.} Las pruebas afirman el \emph{contenido} del mensaje, no
sólo el estado: \code{``3 auto-reparaciones''} y el mensaje del Quality Gate deben aparecer
en \code{blocking\_reason}.

\subsection*{D-01 --- Funciones duplicadas: código muerto con contrato divergente}

\begin{finding}
\textbf{Severidad: media.} Deuda técnica con riesgo de deriva silenciosa.
\end{finding}

\paragraph{Evidencia.} Al preparar las pruebas de fases obsoletas, la cobertura mostró
líneas imposibles de alcanzar. La inspección del archivo reveló que
\code{mark\_downstream\_outdated} y \code{clear\_outdated\_phases} estaban definidas
\textbf{dos veces} cada una en el mismo módulo (líneas 325/351 y 437/478).

\paragraph{Causa raíz.} Python enlaza el nombre a la \emph{última} definición, de modo que el
primer par era código muerto: 34 sentencias que ninguna prueba podía ejecutar. No era
redundancia inocua: los dos contratos diferían. La versión sepultada
\code{mark\_downstream\_outdated} devolvía \code{None} y \emph{sobrescribía} el conjunto de
fases obsoletas; la vigente devuelve la lista acumulada y la \emph{extiende}. Un llamador
que se guiara por la primera firma habría escrito
\code{if mark\_downstream\_outdated(...) is None} y obtenido la respuesta contraria.

\paragraph{Corrección.} Se eliminó el par muerto. El comportamiento observable no cambia
(el nombre ya apuntaba a la segunda definición), y el recuento de sentencias del módulo bajó
de 301 a 267.

\paragraph{Pruebas de regresión.} Una prueba comprueba sobre el código fuente que cada
función está definida una sola vez; otra fija el tipo de retorno (lista) de la función
vigente.

% ===========================================================================
\section{Resultados de ejecución}
\label{sec:resultados}

\subsection*{Antes y después}

Medición sobre el mismo entorno. ``Antes'' corresponde al commit \code{2bddc0a}
(\code{fix: offline mode said it answered the prompt; a badge claimed verification}), el
HEAD inmediatamente anterior a esta campaña, sin ninguno de sus cambios.

\begin{center}
\begin{tabular}{lrrr}
\toprule
\textbf{Métrica} & \textbf{Antes} & \textbf{Después} & \textbf{Variación} \\
\midrule
Pruebas de backend correctas      & 589 & \textbf{825} & +236 \\
Pruebas de backend omitidas       & 3   & 3            & --- \\
Pruebas de frontend correctas     & 64 (10 archivos) & \textbf{73 (11 archivos)} & +9 \\
\midrule
Cobertura global del backend      & 85{,}3\,\% & \textbf{89{,}6\,\%} & +4{,}3 pp \\
Sentencias sin cubrir             & 994 & \textbf{727} & $-267$ \\
\bottomrule
\end{tabular}
\end{center}

\noindent
Dos advertencias de lectura, para no atribuir a esta campaña más de lo que le corresponde:

\begin{itemize}[leftmargin=1.4em,itemsep=1pt]
\item Otra sesión de trabajo entregó cambios en paralelo sobre el mismo árbol (por ejemplo
      en \code{orchestrator/stages/runner.py}), de modo que el movimiento de la cobertura
      \emph{global} no es atribuible sólo a estas pruebas. Las cifras por módulo de la tabla
      siguiente sí lo son: cada una se midió antes y después sobre el mismo módulo.
\item La cifra de ``antes'' de las pruebas de backend es el mismo punto de partida para las
      dos fases; el total combina ambas.
\end{itemize}

\subsection*{Cobertura por módulo priorizado}

\begin{center}
\begin{tabular}{lrrrr}
\toprule
\textbf{Módulo} & \textbf{Antes} & \textbf{Sin cubrir} & \textbf{Después} & \textbf{Sin cubrir} \\
\midrule
\code{services/docker\_service.py}    & 52{,}7\,\% & 69 & \textbf{100{,}0\,\%} & 0 \\
\code{api/routes\_session.py}         & 71{,}3\,\% & 78 & \textbf{100{,}0\,\%} & 0 \\
\code{services/lifecycle\_service.py} & 76{,}4\,\% & 71 & \textbf{100{,}0\,\%} & 0 \\
\code{services/pipeline\_runner.py}   & 82{,}4\,\% & 61 & \textbf{100{,}0\,\%} & 0 \\
\bottomrule
\end{tabular}
\end{center}

\noindent
\code{cost/mlflow\_sink.py} ya estaba al 100\,\% antes de la campaña (33 sentencias); la
corrección de D-03 le añadió 9 sentencias y todas quedan cubiertas (42 de 42).

\subsection*{Frontend: qué se añadió y qué sigue sin cubrirse}

El frontend pasó de 64 a 73 pruebas. Las nuevas no son más pruebas de componente: son
\emph{flujos} que ejecutan los módulos de servicio reales, el cliente HTTP real y las vistas
reales contra un backend simulado en la frontera HTTP (MSW). Esto es lo que faltaba: las 64
pruebas anteriores sustituyen la capa de servicios con \code{vi.mock}, así que ninguna
construía una petición. Un backend que renombrara un parámetro habría mantenido las 64 en
verde.

Lo que ahora queda fijado, y que antes no lo estaba: la URL, el método, los parámetros y el
cuerpo de cada petición del flujo; que la credencial viaja en una cabecera y \emph{no} se
persiste en ningún almacenamiento (Principio VI); que un fallo del backend se propaga en
lugar de inventarse; y --- de extremo a extremo en la interfaz --- que el smoke test informa
``no ejecutado'' cuando el endpoint no responde, en vez de un éxito que nadie observó. Ese
último caso no tenía \emph{ninguna} prueba pese a estar documentado como corregido en
\code{docs/frontend\_audit.md}.

\textbf{Sigue sin cubrirse} el navegador real: no hay Playwright ni Cypress. Estas pruebas
ejercen el flujo con jsdom y un backend simulado, no con un navegador y un backend vivos.

\subsection*{Reproducción}
\label{sec:reproduccion}

\begin{lstlisting}
# Suite de backend con cobertura
.venv/bin/python -m pytest -q \
  --cov=backend/app --cov-report=term-missing:skip-covered

# Suite de frontend
cd frontend && npm test
\end{lstlisting}

\noindent
Evidencia versionada en \code{reports/qe/}:

\begin{itemize}[leftmargin=1.4em,itemsep=1pt]
\item \code{baseline-run.txt} --- salida completa antes (commit \code{2bddc0a}), con la
      tabla de líneas sin cubrir por archivo.
\item \code{after-run.txt} --- salida completa después.
\item \code{coverage-summary.json} --- cobertura por archivo en ambos puntos, en formato
      consumible.
\item \code{hang-trace.txt} --- traza del cuelgue que motivó D-03.
\end{itemize}

% ===========================================================================
\section{Brechas restantes y recomendaciones}
\label{sec:brechas}

Tras esta campaña, los módulos con más sentencias sin cubrir son:

\begin{center}
\begin{tabular}{rrrl}
\toprule
\textbf{Sin cubrir} & \textbf{Cobertura} & \textbf{Sentencias} & \textbf{Módulo} \\
\midrule
67 & 70{,}9\,\% & 230 & \code{services/architecture\_service.py} \\
62 & 63{,}1\,\% & 168 & \code{services/requirements\_service.py} \\
53 & 83{,}3\,\% & 317 & \code{services/model\_sql\_service.py} \\
44 & 70{,}9\,\% & 151 & \code{orchestrator/stages/journal.py} \\
34 & 91{,}6\,\% & 403 & \code{orchestrator/stages/runner.py} \\
29 & 25{,}6\,\% & 39  & \code{ssl\_compat.py} \\
26 & 90{,}3\,\% & 268 & \code{services/security\_service.py} \\
\bottomrule
\end{tabular}
\end{center}

\noindent
\textbf{Recomendación de prioridad para la próxima campaña:}

\begin{enumerate}[leftmargin=1.6em]
\item \code{architecture\_service.py} y \code{model\_sql\_service.py} --- serialización de
      OpenAPI, Mermaid y DDL/JPA. Son artefactos entregables; un error de serialización se
      descubre tarde y en el cliente. Este es el mismo perfil que D-05: código que produce
      archivos que nadie comprueba hasta el final.
\item \code{requirements\_service.py} --- descomposición y refinamiento con LLM.
\item \code{journal.py} --- bordes de presupuesto y serialización del diario de ejecución.
\end{enumerate}

\noindent
\code{ssl\_compat.py} se recomienda \textbf{no} perseguir: es compatibilidad TLS en tiempo de
importación y su cobertura tiene poco valor de producto.

\noindent
El siguiente salto de valor no es más cobertura de línea, sino cerrar la brecha de la
Sección~\ref{sec:limitaciones}: pruebas de aceptación HTTP independientes contra el
microservicio generado y desplegado.

% ===========================================================================
\section{Limitaciones del informe}
\label{sec:limitaciones}

\begin{finding}
\textbf{Cubierto no es correcto.} La cobertura mide qué líneas se ejecutaron, no qué
afirmaciones son verdaderas. Las pruebas fijan el comportamiento \emph{actual} del sistema
frente a los contratos declarados; no demuestran que esos contratos sean los deseados.
\end{finding}

\noindent
Hay una limitación específica y más profunda que conviene que el Project Manager conozca:
\textbf{la verificación que la plataforma hace del código que ella misma genera es
autorreferencial}. El sistema escribe las pruebas Mockito y después las ejecuta; un
\code{mvn test -o} en verde demuestra que el código hace lo que dicen las pruebas de la
plataforma, no lo que dicen los escenarios de aceptación del blueprint. Cerrar esa brecha es
un entregable aparte: pruebas de aceptación HTTP independientes, una por escenario declarado,
contra el servicio desplegado. Queda explícitamente fuera del alcance de esta campaña y se
recomienda como siguiente paso de mayor valor.

\vspace{0.4cm}
\noindent
En el frontend, las 73 pruebas ya no son sólo de componente: 9 de ellas recorren flujos
completos contra un backend simulado en la frontera HTTP. Pero \textbf{no hay navegador
real}: no se usa Playwright ni Cypress, y el backend es un simulador, no un servidor vivo.
Un fallo que sólo aparezca en el navegador (CSS, foco, un evento nativo) no lo detectarían.

\end{document}
